\section{Introduction}
\label{sec:intro}

Single-image metric depth has improved quickly~\cite{bochkovskii2025depthpro,piccinelli2025unidepthv2,
wang2025moge2,hu2024metric3dv2,yang2024dav2}, and these models now run on phones, robots and cars whose
cameras they have never seen. To answer in metres, a model has to resolve a basic ambiguity: the same
pixels fit a wide-angle view of a near scene and a telephoto view of a far one~\cite{he2018focal}.
Some models take the intrinsics as input; others predict them; many simply regress metric depth from
pixels. Whether the latter two groups really infer the camera, or fall back on the size of familiar
objects, is hard to tell from benchmark accuracy, because test sets are dominated by a few cameras.

We propose a test that needs no labels and no new data. Cropping an image and resizing it back is an
\emph{exact} change of camera: the focal length is multiplied by the zoom factor $s$, the principal
point moves in a known way, and every pixel that survives keeps its metric depth (\cref{sec:law}).
A model that understands the camera should therefore keep its depth and scale its focal by $s$.
We ask three questions:
\begin{itemize}
  \item[\textbf{H1}] Do current metric depth models keep the depth of the same scene points under a
        digital zoom, and scale their focal estimate by $s$?
  \item[\textbf{H2}] Do models that receive, predict, or ignore the camera fail in systematically
        different ways, and do these failures follow from a geometric explanation?
  \item[\textbf{H3}] Can the law be used as a training signal on unlabeled images without trading away
        accuracy on ordinary images?
\end{itemize}

\paragraph{Findings.}
For H1, no image-only model passes: Depth Anything V2 shrinks depth by about $1/s$
($\alpha=\DavTwoSAlpha$), and models with a camera head close only part of the gap while
under-scaling the focal (\cref{fig:curves}). For H2, the failures line up with the
\emph{zoom--squash ambiguity} we derive: a zoomed image is identical to an image of the scene compressed
along the optical axis, so cues that read apparent size cannot see a zoom; models that reason this way
sit on the line $\alpha+\beta=1$, and Depth Pro does so per image. Camera-conditioned models are
equivariant with the right intrinsics and inherit exactly $+1$ in $\alpha$ with stale ones. For H3, a
self-consistency objective in the spirit of~\cite{zhong2023equivariance} learns a shortcut, ``blurry
means zoomed''; adding blur invariance removes it and brings $\alpha$ near zero on the training
domains, at a measurable cost on unzoomed and unseen images that we report in full.

\paragraph{Contributions.}
\begin{enumerate}
  \item A label-free \textbf{digital-zoom audit} for metric depth with exact geometry, two exponents for
        depth and focal, and controls for interpolation, detail loss, compression and crop position,
        including a separation of optical from digital zoom (\cref{sec:audit}).
  \item A \textbf{taxonomy of failure modes} over three model groups and a model ladder up to
        Depth Anything~3 and VGGT, with the zoom--squash ambiguity and the $\alpha+\beta=1$ line as a
        falsifiable account of the dominant failure (\cref{sec:law,sec:results}).
  \item An analysis of \textbf{when the label-free score predicts real error}, with scene-level
        uncertainty and per-dataset checks, including where it does not (\cref{sec:results}).
  \item An evaluation of \textbf{label-free correction}: self-consistency fine-tuning, the resolution
        shortcut it learns, EFT+B that removes it, a training-free fix for known zoom factors, and the
        accuracy cost on unseen domains (\cref{sec:eft,sec:eftresults}).
\end{enumerate}
Code, splits and all raw results are released.
